\begin{center}
\LARGE
Solving Equational Constraints via Patterns

\end{center}

\begin{center}
\Large
Olga Caprotti, Hoon Hong

\end{center}

\noindent
Date:  July 19th (Friday)\\
Time:  10:20-10:45\\

\begin{center}
\large
Abstract
\end{center}

Solving equational constraints (over some important domains) is one of core problems in
constraint logic programming. Unluckily, it is also a very difficult one (often uncomputable or
intractable in general). 

In this talk, we present an approach for tackling this problem in which the user is supposed to
provide not only equations but also a ``pattern'', containing unknown variables, of the expected
solutions. We then show how to compute values, if they exist, of the variables that turn the pattern
into a solution. 

It seems that this approach can enlarge the scope of the problems that can be handled. In the talk,
we illustrate this for infinite integral domains, modules and commutative algebras. 

